\documentclass[11pt,a4paper]{article}
\usepackage{amsmath,amssymb,mathtools}
\usepackage{geometry}
\geometry{margin=1in}

\title{\textbf{Note 2: Kitaev's Solution to the Honeycomb Model}}
\date{}

\begin{document}
\maketitle

\section*{1. The Kitaev Honeycomb Model}

Consider a 2D honeycomb lattice with localized spin-1/2 degrees of freedom at each site. The lattice is bipartite, partitioned into sublattices $A$ and $B$, with coordination number $z = 3$.

There are three distinct types of bonds connecting nearest-neighbor $A$ and $B$ sublattice sites:
\begin{enumerate}
    \item $x$-bonds (positive slope)
    \item $y$-bonds (negative slope)
    \item $z$-bonds (vertical)
\end{enumerate}

The spins interact via highly anisotropic nearest-neighbor Ising exchange interactions:
\begin{equation}
    H = -J_x \sum_{x\text{-bonds}} \sigma_j^x \sigma_k^x - J_y \sum_{y\text{-bonds}} \sigma_j^y \sigma_k^y - J_z \sum_{z\text{-bonds}} \sigma_j^z \sigma_k^z,
\end{equation}
where $(J_x, J_y, J_z)$ are coupling constants, and $\sigma_j^\alpha$ ($\alpha \in \{x, y, z\}$) is the Pauli matrix at site $j$.

This is a \textbf{compass model}: the spin component taking part in the exchange interaction is determined via the spatial orientation of the bond. 

Because of this, any given spin cannot minimize all its interactions simultaneously. In other words, $\{\sigma_i^x, \sigma_i^y, \sigma_i^z\}$ do not commute with each other, and thus bond terms do not commute with the Hamiltonian.

\paragraph{Verification of Exchange Frustration:}
Let us define the bond terms sharing site $1$:
\begin{equation}
    K_x = -J_x \sigma_1^x \sigma_2^x, \quad K_y = -J_y \sigma_1^y \sigma_3^y, \quad K_z = -J_z \sigma_1^z \sigma_4^z.
\end{equation}
Computing the commutator of $K_x$ and $K_y$ explicitly:
\begin{align}
    [K_x, K_y] &= J_x J_y [\sigma_1^x \sigma_2^x, \sigma_1^y \sigma_3^y] \nonumber \\
    &= J_x J_y \left( \sigma_1^x \sigma_2^x \sigma_1^y \sigma_3^y - \sigma_1^y \sigma_3^y \sigma_1^x \sigma_2^x \right) \nonumber \\
    &= J_x J_y \left( \sigma_1^x \sigma_1^y - \sigma_1^y \sigma_1^x \right) \sigma_2^x \sigma_3^y \nonumber \\
    &= J_x J_y [\sigma_1^x, \sigma_1^y] \sigma_2^x \sigma_3^y \nonumber \\
    &= 2i J_x J_y \sigma_1^z \sigma_2^x \sigma_3^y \neq 0.
\end{align}
The same applies for all other bond pairs sharing a site. Simultaneous eigenstates are not possible! As described in Note 1, mapping such a spin-1/2 Hamiltonian to free fermionic fields (\textbf{fermionization}) is a standard technique to solve this problem.


\section*{2. Enlarging the Hilbert Space and Fractionalization}

\paragraph{Step 1: Introducing 2 Dirac Fermions per Site}
At a given site right now, we have a 2D Hilbert space, leading to a $2^N \times 2^N$ Hamiltonian space. Replace each spin-1/2 at any site $j$ with 2 spinless Dirac fermions, each of a different flavour (i.e. discernible): $d_{1,j}$ and $d_{2,j}$ (where $d_{\alpha,j}$ is the annihilation operator for flavour-$\alpha$ Dirac fermion at site $j$).

Now, each site is 4D, with occupation number / Fock states given by $\{|0,0\rangle, |1,0\rangle, |0,1\rangle, |1,1\rangle\}$.
For the entire problem, we now have a total Hilbert space dimension of $4^N = 2^{2N}$. This means artificial unphysical states have been introduced at this step, which must be projected out later.

The $d_{1/2,j}$ operators create a complex field, as $d^\dagger \neq d$. We could generate a decomposition of the Pauli matrices at each site into these operators, but we do not.

\paragraph{Step 2: Fractionalize into Majorana Fermions}
A Majorana fermion is defined by a strict reality condition imposed upon its creation/annihilation operators: $\gamma^\dagger = \gamma$, with $\{\gamma^\alpha, \gamma^\beta\} = 2\delta^{\alpha\beta}$. The Majorana fermion is its own antiparticle.

Since we have 2 complex fields $(d_{1,j}, d_{2,j})$ at each site $j$, we extract real and imaginary parts of these fields via obvious linear combinations:
\begin{align}
    \gamma_{1,j} &= \frac{1}{\sqrt{2}} \left( d_{1,j} + d_{1,j}^\dagger \right), \\
    \gamma_{2,j} &= \frac{-i}{\sqrt{2}} \left( d_{1,j} - d_{1,j}^\dagger \right), \\
    \gamma_{3,j} &= \frac{1}{\sqrt{2}} \left( d_{2,j} + d_{2,j}^\dagger \right), \\
    \gamma_{4,j} &= \frac{-i}{\sqrt{2}} \left( d_{2,j} - d_{2,j}^\dagger \right).
\end{align}


\section*{3. Majorana Representation and the $\mathbb{Z}_2$ Gauge Field}

We re-christen these Majorana operators at each site as:
\begin{itemize}
    \item \textbf{Itinerant Majorana (sort of like a matter particle):} $c_j = \gamma_{1,j}$
    \item \textbf{3 Localized Majoranas:} $b_j^x = \gamma_{2,j}, \quad b_j^y = \gamma_{3,j}, \quad b_j^z = \gamma_{4,j}$
\end{itemize}

\paragraph{Verification of Anticommutation Relations of Majorana Fermions:}
Let us verify $\{c_j, b_j^x\}$ explicitly:
\begin{align}
    \{c_j, b_j^x\} &= \left\{ \frac{1}{\sqrt{2}}(d_{1,j} + d_{1,j}^\dagger), \frac{-i}{\sqrt{2}}(d_{1,j} - d_{1,j}^\dagger) \right\} \nonumber \\
    &= \frac{-i}{2} \left( \{d_{1,j}, d_{1,j}\} - \{d_{1,j}, d_{1,j}^\dagger\} + \{d_{1,j}^\dagger, d_{1,j}\} - \{d_{1,j}^\dagger, d_{1,j}^\dagger\} \right) \nonumber \\
    &= \frac{-i}{2} \left( 0 - 1 + 1 - 0 \right) = 0.
\end{align}
Similarly, all different Majorana operators on-site and off-site anti-commute cleanly. Fermionization is achieved as the algebra is no longer hard-core boson-esque!

\paragraph{Step 3: Decompose Pauli Matrices into Majorana Operators}
Kitaev proposed the following decomposition:
\begin{equation}
    \sigma_j^\alpha = i b_j^\alpha c_j \quad \forall\, \alpha \in \{x,y,z\}.
\end{equation}

Let us look at how this changes the Kitaev interaction term:
\begin{align}
    -J_\alpha \sigma_j^\alpha \sigma_k^\alpha &= -J_\alpha (i b_j^\alpha c_j)(i b_k^\alpha c_k) \nonumber \\
    &= -J_\alpha (i^2) (b_j^\alpha c_j b_k^\alpha c_k) \nonumber \\
    &= J_\alpha (i b_j^\alpha b_k^\alpha) c_j c_k.
\end{align}
Why? This isolates the interacting term effectively!

Redefine as a local bond operator:
\begin{equation}
    u_{jk} = i b_j^\alpha b_k^\alpha.
\end{equation}
$u_{jk}$ is defined on the bonds. Since all interaction terms now possess this $u_{jk}$ factor, we have $u_{\langle j,k \rangle} \Rightarrow u_{jk}$ are constants of motion, i.e., $[u_{jk}, H] = 0$. These allow for $c$-Majorana hopping!

Note:
\begin{align}
    (u_{jk})^2 &= (i b_j^\alpha b_k^\alpha)(i b_j^\alpha b_k^\alpha) \nonumber \\
    &= (-1) (b_j^\alpha b_k^\alpha b_j^\alpha b_k^\alpha) \nonumber \\
    &= (-1) (-1) (b_j^\alpha)^2 (b_k^\alpha)^2 = 1.
\end{align}
Eigenvalues of $u_{jk} = \{-1, +1\}$; it is a $\mathbb{Z}_2$ variable!

\paragraph{Verification of the Commutation Relations of $u_{jk}$:}
Let us focus on some site $j$, connected to an $x$-bond $\langle j,i \rangle$ and a $y$-bond $\langle j,k \rangle$:
\begin{align}
    u_{ji} u_{jk} &= (i b_j^x b_i^x)(i b_j^y b_k^y) = - b_j^x b_i^x b_j^y b_k^y = - b_j^x (-b_j^y b_i^x) b_k^y = b_j^x b_j^y b_i^x b_k^y, \\
    u_{jk} u_{ji} &= (i b_j^y b_k^y)(i b_j^x b_i^x) = - b_j^y b_k^y b_j^x b_i^x = - b_j^y (-b_j^x b_k^y) b_i^x = b_j^y b_j^x b_k^y b_i^x = - b_j^x b_j^y b_i^x b_k^y.
\end{align}
Wait! Let us trace the swaps carefully:
$b_k^y$ anticommutes with $b_j^x$ (diff flavour) $\to$ 1 sign.
$b_j^y$ anticommutes with $b_j^x$ (diff flavour) $\to$ 1 sign.
$b_k^y$ anticommutes with $b_i^x$ (diff site) $\to$ 1 sign.
Total 3 anticommutations $\to (-1)^3 = -1$.
Combined with the explicit $-1$ from $i^2$:
\begin{equation}
    u_{ji} u_{jk} = u_{jk} u_{ji} \implies [u_{ji}, u_{jk}] = 0.
\end{equation}
The bond operators commute among themselves!

Now, the Hamiltonian in the Majorana basis takes the quadratic form:
\begin{equation}
    H = \frac{i}{4} \sum_{\langle j,k \rangle} A_{jk} c_j c_k, \quad \text{where } A_{jk} = 2 J_\alpha u_{jk}.
\end{equation}


\section*{4. $[u_{jk}, H] = 0$ and the Physical Projection Operator}

We need to show $[u_{jk}, H] = 0$ for $u_{jk}$ to be constants of motion.
For any term $c_l c_m$ in $H$, since $c$-Majoranas are a different flavour from the localized $b$-Majoranas in $u_{jk}$, each $c$ operator anticommutes past $b_j^\alpha$ and $b_k^\alpha$. Two anticommutations lead to a positive sign:
\begin{equation}
    [u_{jk}, c_l c_m] = 0 \implies [u_{jk}, H] = 0.
\end{equation}
Thus, $u_{jk}$ are constants of motion!

To justify the terminology of $\mathbb{Z}_2$ gauge field, we must first isolate the projection operator at each $j$-site that takes us into the appropriate 2D Hilbert space (physical state) from the 4D Majorana Hilbert space. To do this, we require a physical constraint.

The spin operator identity $\sigma_j^x \sigma_j^y \sigma_j^z = i$ is a good candidate, as it involves all Paulis on a site.
Substituting the Majorana description:
\begin{align}
    (i b_j^x c_j)(i b_j^y c_j)(i b_j^z c_j) &= i \nonumber \\
    -i (b_j^x c_j b_j^y c_j)(b_j^z c_j) &= i \nonumber \\
    -(b_j^x b_j^y c_j^2)(b_j^z c_j) &= 1 \nonumber \\
    b_j^x b_j^y b_j^z c_j &= 1.
\end{align}
This must hold for physical states! We define this as a projection operator on each site:
\begin{equation}
    D_j = b_j^x b_j^y b_j^z c_j.
\end{equation}
Essentially, $|\psi\rangle$ is a valid physical state if it is an eigenstate of $D_j$ with eigenvalue $+1$ for all sites $j$.

Let us first represent $D_j$ in the intermediate 2-flavour Dirac fermion representation:
\begin{align}
    D_j &= b_j^x b_j^y b_j^z c_j \nonumber \\
    &= \left[\frac{-i}{\sqrt{2}}(d_{1,j}-d_{1,j}^\dagger)\right] \left[\frac{1}{\sqrt{2}}(d_{2,j}+d_{2,j}^\dagger)\right] \left[\frac{-i}{\sqrt{2}}(d_{2,j}-d_{2,j}^\dagger)\right] \left[\frac{1}{\sqrt{2}}(d_{1,j}+d_{1,j}^\dagger)\right] \nonumber \\
    &= -\frac{1}{4} (d_{1,j}-d_{1,j}^\dagger)(d_{2,j}+d_{2,j}^\dagger)(d_{2,j}-d_{2,j}^\dagger)(d_{1,j}+d_{1,j}^\dagger) \nonumber \\
    &= (1 - 2d_{1,j}^\dagger d_{1,j})(1 - 2d_{2,j}^\dagger d_{2,j}).
\end{align}

Evaluating $D_j$ on Fock states in $d$-fermion occupation number representation:
\begin{align}
    D_j |0,0\rangle_j &= +1 |0,0\rangle_j \quad \text{(Physical, Even particle sector)} \\
    D_j |1,1\rangle_j &= +1 |1,1\rangle_j \quad \text{(Physical, Even particle sector)} \\
    D_j |1,0\rangle_j &= -1 |1,0\rangle_j \quad \text{(Unphysical, Odd particle sector)} \\
    D_j |0,1\rangle_j &= -1 |0,1\rangle_j \quad \text{(Unphysical, Odd particle sector)}
\end{align}

So, acting with $\frac{1+D_j}{2}$ on this space preserves even sector states ($D_j$ sends them to 1) and zeroes out odd sector states ($D_j$ sends them to $-1$).
The global projector is:
\begin{equation}
    P = \prod_{j} \left( \frac{1 + D_j}{2} \right).
\end{equation}

\paragraph{How does $D_j$ interact with $u_{jk}$?}
\begin{equation}
    D_j u_{jk} = (b_j^x b_j^y b_j^z c_j) (i b_j^\alpha b_k^\alpha).
\end{equation}
- $b_k^\alpha$ (different site): 4 anticommutations with $b_j^x, b_j^y, b_j^z, c_j \to (+1)$.
- $c_j$: 1 anticommulation with $b_j^\alpha \to (-1)$.
- $b_j^\alpha$: commutes with itself, 2 anticommutations with other two $b$-flavours $\to (+1)$.
Net result: $D_j u_{jk} + u_{jk} D_j = 0 \implies \{u_{jk}, D_j\} = 0$.

This means $u_{jk}$ is not gauge-invariant, i.e., it is not a physical observable, and there will be different sets of $u_{jk}$ that lead to the same physical states. They form a $\mathbb{Z}_2$ gauge field.


\section*{5. Bilinear Form and Space Dimension Counting}

Thus, the Hamiltonian can now be written as:
\begin{equation}
    H = \Phi^T H([u]) \Phi,
\end{equation}
where $\Phi = (c_1, c_2, \dots, c_N)^T$, and $H([u])$ is an $N \times N$ matrix that depends only on the values of $u_{jk}$. (Exactly 3 numbers per row in the matrix give Hamiltonian hopping amplitudes). This is a tight-binding Hamiltonian of itinerant Majoranas coupled to $u_{jk}$ on the bonds.

\paragraph{Hilbert Space Dimension Counting:}
How do we build the honeycomb lattice? Take an $A$-$B$ sublattice basis connected by a $z$-bond. Translate by unit vectors $\vec{a}_1 = \sqrt{3}\hat{x}$ and $\vec{a}_2 = -\frac{\sqrt{3}}{2}\hat{x} + \frac{3}{2}\hat{y}$.

Originally, let us assume $N$ spins on the lattice.
Originally we have a $2^N \times 2^N$ matrix (Hamiltonian in Pauli basis) that we need to solve.
In the Majorana basis, we end up with 4 Majoranas per site, and thus a $4^N \times 4^N$ problem.

Look at it from the POV of the gauge field. The subsequent projection $D_j = +1$ at each site halves the state dimension ($4 \to 2$ per site), giving $(2^2)^N / 2^N = 2^N$ physical dimension as expected.

\paragraph{Plaquette Operators (Wilson Loops):}
Let us define an index $p$ for each hexagon plaquette. We define:
\begin{equation}
    W_p = \sigma_1^y \sigma_2^z \sigma_3^x \sigma_4^y \sigma_5^z \sigma_6^x = \prod_{\langle j,k \rangle \in \partial p} u_{jk}.
\end{equation}
$W_p^2 = 1 \implies W_p$ has eigenvalues $\pm 1$.
$[W_p, H] = 0$ and $[W_p, W_q] = 0$. So, $W_p$'s and $u_{ij}$'s form the relevant integrals of motion in the Kitaev model.


\section*{6. Direct Sum Partitioning and Block-Diagonality Proof}

By the spectral theorem, the full extended Hilbert space may be partitioned into a direct sum over simultaneous eigenspaces:
\begin{equation}
    H = \bigoplus_{u} H_u.
\end{equation}
Each subspace $H_u$ is spanned by states sharing the exact same background gauge field $u = \{u_{ij}\}$.

\paragraph{Proof that $H$ is Block Diagonal across Gauge Sectors:}
Let $|\psi_u\rangle$ and $|\psi_{u'}\rangle$ be eigenvectors of link operators $\hat{u}_{ij}$, differing at least at one bond $\langle K,L \rangle$ such that $u_{KL}' = -u_{KL}$:
\begin{equation}
    \hat{u}_{ij} |\psi_u\rangle = u_{ij} |\psi_u\rangle, \quad \hat{u}_{ij} |\psi_{u'}\rangle = u_{ij}' |\psi_{u'}\rangle.
\end{equation}

We evaluate the matrix element $\langle \psi_{u'} | \hat{u}_{KL} H | \psi_u \rangle$:
\begin{align}
    \langle \psi_{u'} | \hat{u}_{KL} H | \psi_u \rangle &= (\hat{u}_{KL} |\psi_{u'}\rangle)^\dagger H |\psi_u\rangle = u_{KL}' \langle \psi_{u'} | H | \psi_u \rangle.
\end{align}
Alternatively, using $[\hat{u}_{KL}, H] = 0$:
\begin{align}
    \langle \psi_{u'} | \hat{u}_{KL} H | \psi_u \rangle &= \langle \psi_{u'} | H \hat{u}_{KL} | \psi_u \rangle = u_{KL} \langle \psi_{u'} | H | \psi_u \rangle.
\end{align}
Equating the two:
\begin{equation}
    u_{KL}' \langle \psi_{u'} | H | \psi_u \rangle = u_{KL} \langle \psi_{u'} | H | \psi_u \rangle \implies (u_{KL}' - u_{KL}) \langle \psi_{u'} | H | \psi_u \rangle = 0.
\end{equation}
Since $u_{KL}' \neq u_{KL}$, we have $u_{KL}' - u_{KL} = -2u_{KL} \neq 0$. Therefore:
\begin{equation}
    \langle \psi_{u'} | H | \psi_u \rangle = 0.
\end{equation}
This proves $H$ is strictly block diagonal!


\section*{7. Lieb's Theorem and Projection to Physical Space}

What flux distribution of $W_p$'s yields the ground state?

\paragraph{Lieb's Theorem:}
For tri-coordinated lattices possessing appropriate mirror symmetries, the flux through each elementary plaquette in the absolute ground state is constrained:
\begin{itemize}
    \item Plaquettes of length $2 \pmod 4$ carry zero flux, or $W_p = +1$.
    \item Plaquettes of length $0 \pmod 4$ carry $\pi$-flux, or $W_p = -1$.
\end{itemize}
The honeycomb plaquette has 6 vertices ($6 \equiv 2 \pmod 4$). Thus, the ground state of the Kitaev honeycomb model must have $W_p = +1$ (vortex-free sector).
An easy choice is to set $u_{ij} = +1$ for all bonds $\langle i,j \rangle$.

\paragraph{Projecting Back to Physical Space:}
Let $|\Psi\rangle_{\text{ext}}$ be a state in the extended space corresponding to the zero-vortex $u_{ij} = +1$ configuration. Then the physical state is derived as:
\begin{equation}
    |\Psi\rangle_{\text{phys}} = P |\Psi\rangle_{\text{ext}}.
\end{equation}


\section*{8. Complete Diagonalization of $A[u]$ via $SO(2N)$}

For a chosen background $u$, the itinerant $c$-Majorana Hamiltonian is:
\begin{equation}
    H = \frac{i}{4} \sum_{j,k} A_{jk}[u] c_j c_k,
\end{equation}
where $A[u]$ is a real, anti-symmetric $2N \times 2N$ matrix ($A^T = -A$).

Since $A$ is real and anti-symmetric, $iA[u]$ is Hermitian. Its non-zero eigenvalues occur in equal and opposite pairs $\pm \epsilon_m$ ($\epsilon_m > 0$).

Let $\Psi_m$ be an eigenvector of $iA[u]$ with eigenvalue $\epsilon_m$:
\begin{equation}
    i A \Psi_m = \epsilon_m \Psi_m \implies A \Psi_m = -i \epsilon_m \Psi_m.
\end{equation}
Decomposing $\Psi_m$ into real and imaginary parts $u_m, v_m \in \mathbb{R}^{2N}$:
\begin{equation}
    \Psi_m = u_m + i v_m.
\end{equation}
Substituting:
\begin{equation}
    A (u_m + i v_m) = -i \epsilon_m (u_m + i v_m) = \epsilon_m v_m - i \epsilon_m u_m.
\end{equation}
Equating real and imaginary parts:
\begin{align}
    A u_m &= \epsilon_m v_m, \\
    A v_m &= -\epsilon_m u_m.
\end{align}

Thus, $A$ acts on the 2D real subspace spanned by $\{u_m, v_m\}$ and maps them into each other.
We construct a real $2N \times 2N$ matrix $Q \in SO(2N)$ using $u_m$ and $v_m$ as columns:
\begin{equation}
    Q = \begin{pmatrix} \sqrt{2} u_1 & \sqrt{2} v_1 & \sqrt{2} u_2 & \sqrt{2} v_2 & \dots & \sqrt{2} u_N & \sqrt{2} v_N \end{pmatrix}.
\end{equation}

Evaluating $Q^T A Q$ in the $m$-th $2 \times 2$ block:
\begin{align}
    \begin{pmatrix} u_m^T A u_m & u_m^T A v_m \\ v_m^T A u_m & v_m^T A v_m \end{pmatrix} &= \begin{pmatrix} u_m^T (\epsilon_m v_m) & u_m^T (-\epsilon_m u_m) \\ v_m^T (\epsilon_m v_m) & v_m^T (-\epsilon_m u_m) \end{pmatrix} \nonumber \\
    &= \begin{pmatrix} 0 & -\epsilon_m \\ \epsilon_m & 0 \end{pmatrix}.
\end{align}
Reordering basis columns brings $A$ to canonical skew-symmetric form:
\begin{equation}
    Q^T A Q = \bigoplus_{m=1}^N \begin{pmatrix} 0 & \epsilon_m \\ -\epsilon_m & 0 \end{pmatrix}.
\end{equation}

Transforming $c_j$ into normal mode Majoranas $f_k = \sum_j c_j Q_{jk}$, the Hamiltonian pairs up as:
\begin{equation}
    H = \frac{i}{2} \sum_{m=1}^N \epsilon_m f_{2m-1} f_{2m}.
\end{equation}

Finally, defining complex Dirac fermions:
\begin{equation}
    a_m = \frac{1}{2} (f_{2m-1} + i f_{2m}),
\end{equation}
the Hamiltonian becomes diagonal in the occupation number basis of $a$-fermions:
\begin{equation}
    H = \sum_{m=1}^N \epsilon_m \left( a_m^\dagger a_m - \frac{1}{2} \right),
\end{equation}
with ground state energy $E_0 = -\frac{1}{2} \sum_{m=1}^N \epsilon_m$.

\end{document}